%!TEX program = xelatex
\documentclass[UTF8,11pt]{ctexart}

\usepackage[a4paper,margin=2.2cm]{geometry}
\usepackage{amsmath,amssymb}
\usepackage{booktabs}
\usepackage{longtable}
\usepackage{tabularx}
\usepackage{array}
\usepackage{enumitem}
\usepackage{xcolor}
\usepackage{hyperref}
\usepackage{fontspec}
\usepackage{fancyhdr}

\IfFontExistsTF{PingFang SC}{\setCJKmainfont{PingFang SC}}{}
\IfFontExistsTF{Arial}{\setmainfont{Arial}}{}

\hypersetup{
  colorlinks=true,
  linkcolor=blue!55!black,
  urlcolor=blue!55!black,
  citecolor=blue!55!black
}

\definecolor{phmblue}{RGB}{0,64,152}
\definecolor{phmgray}{RGB}{71,85,105}
\definecolor{phmlight}{RGB}{239,246,255}
\definecolor{phmline}{RGB}{180,202,232}
\definecolor{phmgreen}{RGB}{22,120,83}
\definecolor{phmorange}{RGB}{194,95,20}

\newcommand{\code}[1]{\texttt{\detokenize{#1}}}
\newcommand{\term}[1]{\textbf{\color{phmblue}#1}}
\newcommand{\warn}[1]{\textbf{\color{phmorange}#1}}
\newcommand{\good}[1]{\textbf{\color{phmgreen}#1}}
\newcolumntype{Y}{>{\raggedright\arraybackslash}X}
\newcolumntype{L}[1]{>{\raggedright\arraybackslash}p{#1}}
\setlength{\emergencystretch}{2em}
\setlength{\headheight}{14pt}
\sloppy

\pagestyle{fancy}
\fancyhf{}
\lhead{轴承 PHM 项目术语与特征说明}
\rhead{\thepage}
\renewcommand{\headrulewidth}{0.4pt}

\setlist[itemize]{leftmargin=1.5em,itemsep=0.2em,topsep=0.25em}
\setlist[enumerate]{leftmargin=1.8em,itemsep=0.25em,topsep=0.25em}
\setlength{\parskip}{0.45em}
\setlength{\parindent}{2em}

\title{\heiti 轴承寿命预测与故障诊断系统\\术语、标签、特征与分析指标详解}
\author{PHM 项目文档}
\date{2026 年 6 月 19 日}

\begin{document}
\maketitle

\begin{abstract}
本文档面向不熟悉轴承故障诊断和 PHM 的读者，解释本项目中出现的核心名词、数据形态、标签构造、FPT 与 3$\sigma$ 原则、RUL、健康状态、早期故障检测、故障类型阶段、特征工程、特征分析指标和实验产物。阅读目标是：即使读者不是本领域人员，也能理解本项目为什么要从振动信号中提取特征，为什么要构造 RUL 和健康状态标签，以及为什么某些特征适合寿命预测、某些特征更适合早期故障检测。
\end{abstract}

\tableofcontents
\newpage

\section{先用一句话理解本项目}

本项目研究的是\term{滚动轴承预测性维护}。轴承在机器中负责支撑旋转部件。它从健康到失效通常不是瞬间发生，而是经历一个逐渐退化的过程：开始时振动很平稳，随后出现微小冲击或频率成分变化，最后振动幅值和冲击明显增大，设备接近失效。

本项目要做两类判断：
\begin{enumerate}
  \item \term{还能用多久？} 这对应 \term{RUL}（Remaining Useful Life，剩余使用寿命）预测。
  \item \term{现在处于什么状态？} 这对应健康状态评估、早期故障检测和故障类型阶段识别。
\end{enumerate}

项目中数据流可以概括为：
\[
\text{原始振动 CSV}
\rightarrow
\text{样本索引}
\rightarrow
\text{窗口/快照级特征}
\rightarrow
\text{RUL/健康状态/早期故障标签}
\rightarrow
\text{任务数据集}
\rightarrow
\text{模型训练与分析}
\]

其中最容易混淆的一点是：\warn{原始数据本身通常只是振动信号，不直接告诉我们每一时刻是 healthy、slight、moderate 还是 severe。} 项目里的许多标签是根据轴承全寿命顺序、HI（健康指标）和 FPT（首次可预测时间）构造出来的伪标签或任务标签。

\section{PHM 是什么}

\term{PHM} 是 Prognostics and Health Management 的缩写，中文常译为\term{故障预测与健康管理}。它不是单纯的“分类器”或“回归器”，而是一整套面向设备运维的技术体系。

\begin{longtable}{L{3.2cm}L{11.2cm}}
\toprule
名词 & 通俗解释 \\
\midrule
\endhead
\term{Diagnosis} 诊断 & 判断设备现在是否已经异常、异常属于哪一类。例如：轴承是否有外圈故障。 \\
\term{Prognosis} 预测 & 判断设备未来还能安全运行多久。例如：还剩 30 分钟、50 个采样周期或 20\% 寿命。 \\
\term{Health Management} 健康管理 & 把诊断和预测结果用于维修计划、备件安排、停机风险控制。 \\
\term{Predictive Maintenance} 预测性维护 & 不等坏了再修，也不固定周期盲目换件，而是根据状态和寿命预测安排维修。 \\
\bottomrule
\end{longtable}

本项目中的 PHM 主要落在两个任务上：
\begin{itemize}
  \item \term{RUL 回归}：输出一个连续数值，表示剩余寿命。
  \item \term{状态/故障分类}：输出离散类别，例如 healthy / slight / moderate / severe，或者 early fault 的 0/1。
\end{itemize}

\section{轴承和振动信号的基本概念}

\subsection{轴承为什么能从振动里看出问题}

轴承由内圈、外圈、滚动体和保持架组成。理想情况下，旋转很平稳；当某个部件出现点蚀、裂纹、剥落或磨损时，滚动体经过缺陷位置会产生冲击，冲击会传到传感器上，表现为加速度信号的变化。

\begin{longtable}{L{3.5cm}L{10.8cm}}
\toprule
术语 & 含义 \\
\midrule
\endhead
内圈 Inner Race & 与转轴一起旋转的轴承圈。内圈故障常随旋转调制，特征频率复杂。 \\
外圈 Outer Race & 通常固定在轴承座上。外圈缺陷经过滚动体时会产生周期性冲击。 \\
滚动体 Ball/Roller & 在内圈和外圈之间滚动的球或滚子。滚动体故障可能导致宽频冲击。 \\
保持架 Cage & 分隔滚动体的结构。保持架故障可能引起低频调制或不稳定冲击。 \\
振动加速度 & 本项目主要使用的信号类型，单位可理解为“振动强弱随时间变化的记录”。 \\
水平/垂直通道 & 两个传感器方向。XJTU-SY 和 PHM2012 都有水平、垂直两个振动方向。 \\
\bottomrule
\end{longtable}

\subsection{一个 CSV 文件在项目里代表什么}

本项目把一个原始振动 CSV 文件视为一个\term{振动快照}或\term{样本}。这个样本不是一个单点，而是一小段连续振动信号。

\begin{itemize}
  \item XJTU-SY：每个 CSV 通常是一次采集快照，包含水平和垂直两列振动信号。
  \item PHM2012：每个 \code{acc_*.csv} 通常包含时间列和水平/垂直加速度列；当前框架只保留最后两列加速度。
\end{itemize}

项目中的样本级索引字段包括：
\begin{longtable}{L{3.8cm}L{10.5cm}}
\toprule
字段 & 解释 \\
\midrule
\endhead
\code{sample_uid} & 样本唯一 ID。后续 features、labels、task dataset 都靠它对齐。 \\
\code{dataset} & 数据集名称，例如 \code{XJTU-SY} 或 \code{PHM2012}。 \\
\code{bearing_id} & 轴承 ID，例如 \code{Bearing1_1}。同一个轴承内部样本必须按时间顺序处理。 \\
\code{condition_id} & 工况 ID，例如转速和载荷组合 \code{35Hz12kN}。 \\
\code{source_group} & PHM2012 中区分 \code{Learning_set} 与 \code{Full_Test_Set}；XJTU-SY 通常为空。 \\
\code{sample_id} & 原始文件编号或快照编号。 \\
\code{timestep} & 项目内部的时间顺序编号。它不一定是秒，但保证同一轴承内从早到晚递增。 \\
\bottomrule
\end{longtable}

\section{数据集术语}

\subsection{XJTU-SY}

XJTU-SY 是一个滚动轴承加速退化全寿命数据集，包含 15 个轴承，分为 3 个工况，每个工况 5 个轴承。项目中常用的划分方式是：同一工况下 5 个轴承，每次留 1 个轴承测试，其余轴承用于训练/验证。

\begin{longtable}{L{3.3cm}L{11cm}}
\toprule
名词 & 项目中的含义 \\
\midrule
\endhead
\code{35Hz12kN} & 转速约 35 Hz、载荷 12 kN 的工况。 \\
\code{Bearing1_1}--\code{Bearing1_5} & 该工况下的 5 个轴承。 \\
\code{fault_element} & XJTU-SY 官方轴承级故障元素说明，例如 Outer Race Fault、Cage Fault。它不是逐样本标签。 \\
\code{xjtu_leave_one_bearing_out} & 同一工况下留一个轴承做测试的划分方式。 \\
\bottomrule
\end{longtable}

\subsection{PHM2012 / PRONOSTIA}

PHM2012 是 IEEE PHM 2012 Challenge 使用的 PRONOSTIA 轴承全寿命数据。它同样是 run-to-failure 数据，但不适合强行做故障类型识别，因为没有像 XJTU-SY 那样的官方轴承级 fault element 标签。

\begin{longtable}{L{3.3cm}L{11cm}}
\toprule
名词 & 项目中的含义 \\
\midrule
\endhead
\code{Learning_set} & PHM2012 官方训练/学习集合。 \\
\code{Full_Test_Set} & PHM2012 官方测试/验证集合。 \\
\code{acc_*.csv} & 加速度文件。当前框架读取最后两列作为水平、垂直振动。 \\
\code{temp_*.csv} & 温度文件。当前框架暂不纳入核心特征。 \\
\code{phm2012_official} & 优先按官方 train/test 结构划分。 \\
\bottomrule
\end{longtable}

\section{标签、任务和目标变量}

\subsection{标签 label 是什么}

在机器学习里，\term{标签}是模型要学习预测的目标。输入是特征，输出是标签。例如：
\[
\text{特征} = [\text{RMS}, \text{kurtosis}, \text{spectral entropy}, \ldots]
\quad\Rightarrow\quad
\text{标签} = \text{RUL 或健康状态}
\]

轴承数据里，原始 CSV 通常没有每个快照的精确健康标签。因此项目用以下信息构造标签：
\begin{itemize}
  \item 同一轴承内部的时间顺序。
  \item 轴承最终失效这一事实。
  \item 由特征构造的健康指标 HI。
  \item 由 HI 计算出的 FPT。
  \item XJTU-SY 的轴承级 fault element。
\end{itemize}

\subsection{RUL：Remaining Useful Life}

\term{RUL} 是 Remaining Useful Life，中文是\term{剩余使用寿命}。它回答：\term{从当前时刻到失效还剩多久？}

项目中有三种相关字段：
\begin{longtable}{L{4cm}L{10.2cm}}
\toprule
字段 & 解释 \\
\midrule
\endhead
\code{linear_rul_steps} & 线性 RUL 的剩余样本步数。越接近失效越小。 \\
\code{linear_rul_seconds} & 将剩余步数乘以样本间隔后的秒数。 \\
\code{linear_rul_norm} & 归一化到 $[0,1]$ 的线性 RUL，开始时接近 1，失效时为 0。 \\
\code{piecewise_rul_steps} & 分段 RUL 的剩余步数。FPT 前保持平台，FPT 后下降。 \\
\code{piecewise_rul_seconds} & 分段 RUL 的秒数版本。 \\
\code{piecewise_rul_norm} & 归一化到 $[0,1]$ 的分段 RUL，是 Stage 6 默认分析目标。 \\
\bottomrule
\end{longtable}

\subsubsection{线性 RUL}

如果一个轴承一共有 $N$ 个样本，第 $t$ 个样本的剩余步数为：
\[
\text{linear\_rul\_steps}(t)=N-1-t
\]
归一化后：
\[
\text{linear\_rul\_norm}(t)=\frac{N-1-t}{N-1}
\]

直观理解：从第一天开始就认为寿命在均匀下降。这很简单，但有时不符合物理直觉，因为健康早期的振动可能长期稳定。

\subsubsection{分段 RUL}

分段 RUL 引入 FPT。FPT 前认为设备仍处于健康或稳定阶段，RUL 不明显下降；FPT 后才进入明显退化阶段：
\[
\text{piecewise\_rul\_norm}(t)=
\begin{cases}
1, & t < t_{\mathrm{FPT}} \\
\frac{N-1-t}{N-1-t_{\mathrm{FPT}}}, & t \ge t_{\mathrm{FPT}}
\end{cases}
\]

这更接近项目的直觉：轴承不是从第一分钟就线性坏掉，而是先稳定运行，之后进入可观测退化。

\subsection{退化 degradation}

\term{退化}指设备健康程度随时间变差的过程。它不是某个单一标签，而是一个物理过程。退化通常表现为：
\begin{itemize}
  \item 振动能量上升，例如 RMS、std、mean abs 变大。
  \item 冲击增强，例如 kurtosis、crest factor、impulse factor 增大。
  \item 频谱结构变化，例如 spectral centroid 或 entropy 改变。
  \item 不同健康阶段之间的分布逐渐错开。
\end{itemize}

项目通过 HI、FPT、RUL、健康状态、早期故障等多个对象来描述退化。

\subsection{HI：Health Indicator}

\term{HI} 是 Health Indicator，中文可称为\term{健康指标}。它把一个或多个特征转换成一条随时间变化的健康曲线，用来描述“轴承坏到什么程度”。

当前项目的 \code{degradation_basic} 使用 \term{feature-column HI}：
\begin{enumerate}
  \item 优先选择 \code{mag__time__rms}，如果不存在则选择 \code{h__time__rms}，再不存在则选择 \code{v__time__rms}。
  \item 如果 \code{direction=bad_high}，数值越大表示越坏；如果 \code{bad_low}，会取相反数。
  \item 对同一轴承内的 HI 做滑动平均平滑，默认窗口为 3。
  \item 对每个轴承单独做 min-max 归一化，得到 \code{hi_norm}。
\end{enumerate}

用公式表示，若原始 HI 为 $h_t$，平滑窗口为 $w=3$：
\[
\tilde{h}_t = \frac{1}{|\mathcal{W}_t|}\sum_{i\in\mathcal{W}_t}h_i
\]
归一化为：
\[
h^{\mathrm{norm}}_t=\frac{\tilde{h}_t-\min(\tilde{h})}{\max(\tilde{h})-\min(\tilde{h})}
\]

\warn{重要注意：如果某个特征被用作 HI source，那么它在 FPT、健康状态、早期故障标签上表现好，存在循环定义风险。项目会在 \code{feature_ranking.csv} 和 \code{feature_cards.csv} 中标注 \code{is_label_source=true}。}

\subsection{FPT：First Predictable Time}

\term{FPT} 是 First Predictable Time，中文可理解为\term{首次可预测时间}或\term{首次明显退化点}。它不是最终失效点，而是设备从“看起来健康”进入“退化已经可观测”的分界点。

通俗理解：
\begin{itemize}
  \item FPT 前：信号基本像健康状态，模型很难可靠预测退化。
  \item FPT 后：退化趋势开始显现，RUL 预测和状态评估更有意义。
\end{itemize}

\subsection{3$\sigma$ 原则 / Three Sigma}

当前项目用 3$\sigma$ 原则检测 FPT。做法如下：
\begin{enumerate}
  \item 取每个轴承前 \code{healthy_ratio=0.2} 的样本作为健康基线。
  \item 计算健康基线 HI 的均值 $\mu$ 和标准差 $\sigma$。
  \item 阈值为：
  \[
  \theta = \mu + 3\sigma + \Delta_{\min}
  \]
  当前默认 \code{sigma_ratio=3.0}，\code{min_delta=0.0}。
  \item 从前到后找连续 \code{consecutive_points=3} 个点都超过阈值的位置。
  \item 第一个满足条件的位置就是 FPT。
  \item 如果找不到，使用 fallback。当前默认 fallback 是 \code{healthy_ratio}，即健康段结束处。
\end{enumerate}

为什么叫 3$\sigma$？如果健康期波动大致服从稳定分布，均值加 3 倍标准差以上通常属于很不寻常的偏离。连续多个点越过阈值，比单个点越过更稳健，可以减少偶然尖峰造成误报。

\subsection{健康状态 Health State}

健康状态是把轴承寿命阶段离散化。当前项目默认四类：
\begin{longtable}{L{3cm}L{4cm}L{7cm}}
\toprule
\code{health_state_id} & 名称 & 含义 \\
\midrule
\endhead
0 & \code{healthy} & FPT 前，认为处于健康或稳定阶段。 \\
1 & \code{slight} & FPT 后早期退化，退化刚开始。 \\
2 & \code{moderate} & FPT 后中期退化，趋势更明显。 \\
3 & \code{severe} & FPT 后晚期退化，接近失效。 \\
\bottomrule
\end{longtable}

项目中，FPT 后的剩余退化过程按照进度比例切分。默认边界为 0.4 和 0.8：
\[
\text{progress}=\frac{t-t_{\mathrm{FPT}}}{N-1-t_{\mathrm{FPT}}}
\]
\begin{itemize}
  \item $t<t_{\mathrm{FPT}}$：healthy。
  \item $0\le \text{progress}<0.4$：slight。
  \item $0.4\le \text{progress}<0.8$：moderate。
  \item $\text{progress}\ge 0.8$：severe。
\end{itemize}

\subsection{早期故障 Early Fault}

\term{早期故障检测}在当前项目里是一个二分类标签：
\[
\text{early\_fault}(t)=
\begin{cases}
0, & t<t_{\mathrm{FPT}} \\
1, & t\ge t_{\mathrm{FPT}}
\end{cases}
\]

它回答的是：\term{当前样本是否已经进入 FPT 后的异常/退化阶段？}

注意它不是“故障已经严重”的意思。它更像是早期报警：一旦过了 FPT，就认为需要重点监控。

\subsection{故障类型阶段 Fault Type Stage}

\term{fault type stage} 是项目为 XJTU-SY 做的弱故障类型阶段标签。它结合健康状态和轴承级 fault element：
\begin{itemize}
  \item 健康状态为 0：标为 \code{normal}。
  \item 中间退化阶段：标为 \code{degraded_unknown}。
  \item severe 阶段：使用 XJTU-SY 的轴承级 fault element，例如 outer race、inner race、cage 等。
\end{itemize}

\warn{PHM2012 不应强做 fault type，因为它没有同等的官方 fault element 标签。} 当前分析配置中如果 PHM2012 缺少 \code{fault_type_stage_name}，会在 \code{analysis_report.json} 中记录 \code{fault_type_skipped=true}。

\section{数据划分和泄漏}

\subsection{Train / Validation / Test}

\begin{longtable}{L{3cm}L{11.2cm}}
\toprule
集合 & 作用 \\
\midrule
\endhead
Train 训练集 & 用来拟合模型、估计标准化参数、做特征 ranking。 \\
Validation 验证集 & 用来调参、早停、选择模型。 \\
Test 测试集 & 只用于最终评估，不能用于选择特征或调参。 \\
\bottomrule
\end{longtable}

本项目特别强调：\term{特征选择、标准化参数、特征 ranking 应尽量只用 train split。} 如果提前看 test label 来挑特征，结果会显得很好，但泛化评估不可信。

\subsection{数据泄漏 Data Leakage}

\term{数据泄漏}指模型或分析过程使用了本不应该提前知道的信息。例如：
\begin{itemize}
  \item 用测试集均值和标准差标准化训练集。
  \item 用测试集标签挑选最好的特征。
  \item 用某个特征构造标签，又宣称同一个特征独立发现了该标签。
\end{itemize}

本项目最需要标注的是 \term{label-source leakage}。例如 \code{mag__time__rms} 被用作 HI source，FPT 和 early fault 标签又由 HI 构造，那么 \code{mag__time__rms} 在 early fault 上表现好并不意外。它仍可用，但论文和报告中必须说明。

\section{原始信号、通道和特征命名}

\subsection{通道命名}

当前 manual 特征使用三类通道：
\begin{longtable}{L{2.5cm}L{11.7cm}}
\toprule
通道 & 含义 \\
\midrule
\endhead
\code{h} & horizontal，水平方向振动。 \\
\code{v} & vertical，垂直方向振动。 \\
\code{mag} & magnitude，水平和垂直的合成幅值，计算为 $\sqrt{h^2+v^2}$。 \\
\bottomrule
\end{longtable}

\subsection{特征列命名规则}

manual 特征列遵循：
\[
\code{channel\_\_domain\_\_feature}
\]
例如：
\begin{itemize}
  \item \code{h__time__rms}：水平通道的时间域 RMS。
  \item \code{v__spectral__entropy}：垂直通道的频域谱熵。
  \item \code{mag__time__kurtosis}：合成幅值通道的时间域峭度。
\end{itemize}

tsfresh 特征列遵循：
\[
\code{tsfresh\_\_kind\_\_calculator\_\_params}
\]
其中 \code{kind} 通常是 \code{h} 或 \code{v}，例如某个通道的均值、方差、FFT 系数、自相关等。

\section{时间域特征}

时间域特征直接在原始振动序列 $x_1,x_2,\ldots,x_N$ 上计算。它们容易理解，适合向非专业读者解释。

\subsection{manual\_basic 当前启用的时间域特征}

令 $x_i$ 是一个快照内第 $i$ 个振动采样点，$N$ 是该快照采样点数，$|x_i|$ 是绝对值。

\begin{longtable}{L{3.5cm}L{4.2cm}L{6.5cm}}
\toprule
特征名 & 公式/定义 & 非专业解释与用途 \\
\midrule
\endhead
\code{mean} & $\bar{x}=\frac{1}{N}\sum_i x_i$ & 平均振动偏置。若传感器零点偏移或信号不对称，mean 会变化。通常不是最强故障特征。 \\
\code{std} & $\sqrt{\frac{1}{N}\sum_i(x_i-\bar{x})^2}$ & 振动离散程度。退化加剧时，振动波动通常更大，适合 RUL 和健康状态。 \\
\code{rms} & $\sqrt{\frac{1}{N}\sum_i x_i^2}$ & 均方根，代表振动能量强度。轴承越坏，RMS 常越高，是最常用退化特征之一。 \\
\code{mean_abs} & $\frac{1}{N}\sum_i |x_i|$ & 平均绝对振动幅值。比 mean 更能反映振动强度，适合 RUL 和健康状态。 \\
\code{ptp} & $\max(x)-\min(x)$ & peak-to-peak，峰峰值。反映一个快照内最大振动跨度，对冲击和严重退化敏感。 \\
\code{skewness} & 三阶标准化矩 & 偏度，描述波形是否向某一侧偏。可反映不对称冲击，但解释性弱于 RMS。 \\
\code{kurtosis} & 四阶标准化矩 & 峭度，描述尖峰/重尾程度。对局部冲击非常敏感，常用于早期故障检测。 \\
\code{crest_factor} & $\frac{\max |x_i|}{\mathrm{RMS}}$ & 峰值因子。看“最高尖峰”相对整体能量有多突出，适合发现突然冲击。 \\
\code{impulse_factor} & $\frac{\max |x_i|}{\mathrm{mean\_abs}}$ & 脉冲因子。类似 crest factor，但以平均绝对幅值为基准，对尖峰敏感。 \\
\code{clearance_factor} & $\frac{\max |x_i|}{(\frac{1}{N}\sum_i\sqrt{|x_i|})^2}$ & 裕度因子。常用于滚动轴承冲击类故障，对尖锐冲击敏感。 \\
\bottomrule
\end{longtable}

\subsection{代码支持但 manual\_basic 未默认启用的时间域特征}

\begin{longtable}{L{3.5cm}L{10.7cm}}
\toprule
特征名 & 解释 \\
\midrule
\endhead
\code{var} & 方差，等于 std 的平方，表示振动波动能量。 \\
\code{max} & 最大值，可能受单个异常点影响。 \\
\code{min} & 最小值，常与 max/ptp 配合理解。 \\
\code{shape_factor} & $\mathrm{RMS}/\mathrm{mean\_abs}$，形状因子，反映波形形态。 \\
\bottomrule
\end{longtable}

\section{频域特征}

频域特征先把时域信号通过 rFFT 转成频谱。频谱可以理解为：振动能量分布在哪些频率上。

对每个快照，项目计算一侧实数 FFT：
\[
X_k = \mathrm{rFFT}(x)_k,\quad P_k=|X_k|^2
\]
其中 $P_k$ 是第 $k$ 个频率 bin 的功率。默认 \code{include_dc=false}，也就是去掉 0 Hz 的直流分量。

\subsection{manual\_basic 当前启用的频域特征}

\begin{longtable}{L{3.8cm}L{4.4cm}L{6.1cm}}
\toprule
特征名 & 公式/定义 & 非专业解释与用途 \\
\midrule
\endhead
\code{centroid} & $\sum_k f_k p_k$ & 频谱质心。把频谱看作质量分布，centroid 是能量中心频率。退化可能导致能量向高频或低频迁移。 \\
\code{bandwidth} & $\sqrt{\sum_k(f_k-c)^2p_k}$ & 频谱带宽。表示能量围绕质心分散得多宽。冲击和宽频噪声会影响它。 \\
\code{rms_frequency} & $\sqrt{\sum_k f_k^2p_k}$ & RMS 频率。表示频率成分的均方根位置，常用于频率结构变化分析。 \\
\code{peak_frequency} & $\arg\max_{f_k} P_k$ & 功率最大的频率。看当前最强的振动频率在哪里。 \\
\code{entropy} & $-\sum_k p_k\log(p_k)/\log K$ & 谱熵。频谱越分散、越复杂，熵越高；频谱集中在少数频率时熵低。 \\
\bottomrule
\end{longtable}

这里 $p_k=P_k/\sum_j P_j$ 是归一化频谱功率，$c$ 是 centroid。

\subsection{代码支持但 manual\_basic 未默认启用的频域特征}

\begin{longtable}{L{3.5cm}L{10.7cm}}
\toprule
特征名 & 解释 \\
\midrule
\endhead
\code{flatness} & 谱平坦度。频谱越像白噪声，flatness 越高；越像单一尖峰，flatness 越低。 \\
\code{rolloff} & 谱滚降频率。累计能量达到某个比例（默认 85\%）时对应的频率。 \\
\bottomrule
\end{longtable}

\section{manual\_basic 的完整特征列}

\code{manual_basic} 当前启用 3 个通道，每个通道 10 个时间域特征和 5 个频域特征，因此一共：
\[
3\times(10+5)=45
\]
个 manual 特征。

\subsection{时间域 30 个特征列}

\begin{longtable}{L{4.4cm}L{9.7cm}}
\toprule
通道 & 特征列 \\
\midrule
\endhead
\code{h} & \code{h__time__mean}, \code{h__time__std}, \code{h__time__rms}, \code{h__time__mean_abs}, \code{h__time__ptp}, \code{h__time__skewness}, \code{h__time__kurtosis}, \code{h__time__crest_factor}, \code{h__time__impulse_factor}, \code{h__time__clearance_factor} \\
\code{v} & \code{v__time__mean}, \code{v__time__std}, \code{v__time__rms}, \code{v__time__mean_abs}, \code{v__time__ptp}, \code{v__time__skewness}, \code{v__time__kurtosis}, \code{v__time__crest_factor}, \code{v__time__impulse_factor}, \code{v__time__clearance_factor} \\
\code{mag} & \code{mag__time__mean}, \code{mag__time__std}, \code{mag__time__rms}, \code{mag__time__mean_abs}, \code{mag__time__ptp}, \code{mag__time__skewness}, \code{mag__time__kurtosis}, \code{mag__time__crest_factor}, \code{mag__time__impulse_factor}, \code{mag__time__clearance_factor} \\
\bottomrule
\end{longtable}

\subsection{频域 15 个特征列}

\begin{longtable}{L{4.4cm}L{9.7cm}}
\toprule
通道 & 特征列 \\
\midrule
\endhead
\code{h} & \code{h__spectral__centroid}, \code{h__spectral__bandwidth}, \code{h__spectral__rms_frequency}, \code{h__spectral__peak_frequency}, \code{h__spectral__entropy} \\
\code{v} & \code{v__spectral__centroid}, \code{v__spectral__bandwidth}, \code{v__spectral__rms_frequency}, \code{v__spectral__peak_frequency}, \code{v__spectral__entropy} \\
\code{mag} & \code{mag__spectral__centroid}, \code{mag__spectral__bandwidth}, \code{mag__spectral__rms_frequency}, \code{mag__spectral__peak_frequency}, \code{mag__spectral__entropy} \\
\bottomrule
\end{longtable}

\section{tsfresh 特征}

\term{tsfresh} 是一个自动时间序列特征提取库。它会从时间序列中计算大量统计、频域、复杂度和自相关特征。项目提供：
\begin{itemize}
  \item \code{tsfresh_minimal}：最小特征集合，速度快、解释简单。
  \item \code{tsfresh_efficient}：更丰富的特征集合，覆盖更多统计特征，但列数多、解释成本更高。
  \item \code{manual_tsfresh_basic}：manual 特征 + tsfresh minimal 的组合。
\end{itemize}

\subsection{tsfresh minimal 的特征计算器}

当前环境中 \code{MinimalFCParameters} 包含：
\begin{longtable}{L{4.2cm}L{10cm}}
\toprule
计算器 & 通俗解释 \\
\midrule
\endhead
\code{sum_values} & 所有采样点求和。 \\
\code{median} & 中位数。 \\
\code{mean} & 均值。 \\
\code{length} & 序列长度。 \\
\code{standard_deviation} & 标准差。 \\
\code{variance} & 方差。 \\
\code{root_mean_square} & RMS。 \\
\code{maximum} & 最大值。 \\
\code{absolute_maximum} & 绝对值最大值。 \\
\code{minimum} & 最小值。 \\
\bottomrule
\end{longtable}

\subsection{tsfresh efficient 的特征类型}

\code{EfficientFCParameters} 当前包含 73 类计算器，其中许多计算器会展开成多个参数列。它覆盖：
\begin{itemize}
  \item 基本统计：mean、median、variance、standard deviation、skewness、kurtosis。
  \item 能量和变化：abs energy、mean abs change、absolute sum of changes。
  \item 重复和计数：duplicate、value count、range count、count above/below mean。
  \item 自相关和局部相关：autocorrelation、partial autocorrelation、agg autocorrelation。
  \item 频域：FFT coefficient、FFT aggregated、Welch density、fourier entropy。
  \item 小波和峰值：cwt coefficients、number of peaks、number cwt peaks。
  \item 趋势和稳定性：linear trend、augmented Dickey-Fuller、large standard deviation。
  \item 复杂度和信息量：permutation entropy、Lempel-Ziv complexity、binned entropy。
\end{itemize}

非专业读者可以这样理解：manual 特征像是“人工挑选的少数可解释指标”；tsfresh 像是“自动生成的候选指标库”。tsfresh 可能找到有用信号，但每个特征的物理意义通常不如 RMS、kurtosis、spectral entropy 直观。

\subsection{tsfresh efficient 当前计算器名称总览}

下面列出项目当前环境中 \code{EfficientFCParameters} 的计算器名称。这里的“计算器”不是最终列名；很多计算器会因为不同参数组合展开成多列。例如 \code{autocorrelation} 会按多个 lag 展开，\code{fft_coefficient} 会按不同 coeff 和 attr 展开。

\begin{longtable}{L{4.2cm}L{10cm}}
\toprule
类别 & 计算器名称 \\
\midrule
\endhead
基础统计 & \code{sum_values}, \code{median}, \code{mean}, \code{length}, \code{standard_deviation}, \code{variance}, \code{variation_coefficient}, \code{maximum}, \code{minimum}, \code{absolute_maximum}, \code{root_mean_square}, \code{skewness}, \code{kurtosis} \\
重复与计数 & \code{has_duplicate}, \code{has_duplicate_max}, \code{has_duplicate_min}, \code{value_count}, \code{range_count}, \code{count_above}, \code{count_below}, \code{count_above_mean}, \code{count_below_mean}, \code{ratio_value_number_to_time_series_length} \\
极值位置 & \code{first_location_of_maximum}, \code{last_location_of_maximum}, \code{first_location_of_minimum}, \code{last_location_of_minimum}, \code{longest_strike_above_mean}, \code{longest_strike_below_mean} \\
变化与能量 & \code{abs_energy}, \code{absolute_sum_of_changes}, \code{mean_abs_change}, \code{mean_change}, \code{mean_second_derivative_central}, \code{cid_ce}, \code{number_crossing_m}, \code{number_peaks}, \code{number_cwt_peaks}, \code{mean_n_absolute_max} \\
分布与分位数 & \code{quantile}, \code{index_mass_quantile}, \code{large_standard_deviation}, \code{ratio_beyond_r_sigma}, \code{variance_larger_than_standard_deviation}, \code{symmetry_looking}, \code{benford_correlation} \\
重复值统计 & \code{percentage_of_reoccurring_values_to_all_values}, \code{percentage_of_reoccurring_datapoints_to_all_datapoints}, \code{sum_of_reoccurring_values}, \code{sum_of_reoccurring_data_points} \\
相关与自回归 & \code{autocorrelation}, \code{agg_autocorrelation}, \code{partial_autocorrelation}, \code{ar_coefficient}, \code{c3}, \code{time_reversal_asymmetry_statistic} \\
频域与谱特征 & \code{fft_coefficient}, \code{fft_aggregated}, \code{spkt_welch_density}, \code{fourier_entropy} \\
小波与局部形态 & \code{cwt_coefficients}, \code{change_quantiles} \\
趋势与平稳性 & \code{linear_trend}, \code{agg_linear_trend}, \code{linear_trend_timewise}, \code{augmented_dickey_fuller} \\
复杂度与信息量 & \code{binned_entropy}, \code{permutation_entropy}, \code{lempel_ziv_complexity}, \code{energy_ratio_by_chunks}, \code{query_similarity_count} \\
非线性动力学 & \code{friedrich_coefficients}, \code{max_langevin_fixed_point} \\
\bottomrule
\end{longtable}

如果后续使用 \code{tsfresh_efficient}，建议先用 Stage 6 的 \code{feature_summary.csv} 删除常数列和异常列，再用 \code{feature_ranking.csv} 和 \code{feature_cards.csv} 选择少量稳定、可解释的特征。不要直接把上千列 tsfresh 特征全部放进论文解释中。

\section{特征清洗和标准化}

特征提取后，项目会进入 \term{FeatureCleaner}：
\begin{enumerate}
  \item \term{median impute}：用训练集特征中位数填补 NaN/Inf。
  \item \term{standard scale}：标准化为均值 0、标准差 1。
  \item \term{drop constant}：删除训练集中几乎不变的常数特征。
  \item \term{train-only fit}：中位数、均值、标准差只从训练集估计。
\end{enumerate}

标准化公式：
\[
z=\frac{x-\mu_{\mathrm{train}}}{\sigma_{\mathrm{train}}}
\]

\warn{标准化参数不能从测试集估计。} 否则模型在训练阶段就“看见”了测试集分布，会导致评估偏乐观。

\section{特征分析 Stage 6}

Stage 6 的目标不是训练模型，而是回答：
\begin{itemize}
  \item 哪些特征适合 RUL？
  \item 哪些特征适合健康状态？
  \item 哪些特征适合早期故障检测？
  \item 哪些特征只是末期突变，不适合早期检测？
  \item 哪些特征参与了 HI/FPT 标签构造，存在 label-source 风险？
\end{itemize}

Stage 6 的主要产物：
\begin{longtable}{L{5cm}L{9.2cm}}
\toprule
文件 & 含义 \\
\midrule
\endhead
\code{feature_summary.csv} & 每个特征的缺失、方差、均值、分位数、常数列检查。 \\
\code{rul_correlation.csv} & 特征与 RUL 的 Pearson/Spearman/Kendall 相关性。 \\
\code{degradation_scores.csv} & 单调性、鲁棒性、趋势一致性、末端一致性。 \\
\code{health_state_separability.csv} & 健康状态可分性，包括 MI、Fisher score、ANOVA F 和各状态均值/标准差。 \\
\code{early_fault_scores.csv} & 早期故障敏感性，包括 Cohen's d、AUC、FPT 前后均值差。 \\
\code{fault_type_scores.csv} & 故障类型阶段的弱可分性指标。 \\
\code{feature_ranking.csv} & 综合 ranking，包含 RUL/Health/Early/FaultType 分数和推荐任务。 \\
\code{feature_cards.csv} & 每个特征的解释卡片：物理意义、推荐任务、原因、 caveat、泄漏标记。 \\
\code{feature_recommendations.md} & 面向读者的文字推荐报告。 \\
\code{figures/} & 特征曲线、箱线图、热图和任务推荐矩阵。 \\
\bottomrule
\end{longtable}

\section{RUL 特征分析指标}

RUL 关注：\term{这个特征能不能稳定反映退化趋势？}

\subsection{相关性}

项目计算 Pearson、Spearman、Kendall：
\begin{longtable}{L{3.5cm}L{10.7cm}}
\toprule
指标 & 解释 \\
\midrule
\endhead
Pearson & 线性相关。适合关系近似直线的特征。 \\
Spearman & 排序相关。只要一个特征随退化大致单调变化，即使不是直线，也能反映出来。 \\
Kendall & 排序一致性，更保守。 \\
\bottomrule
\end{longtable}

因为 RUL 越接近失效越小，而许多退化特征越接近失效越大，所以常见的是负相关。例如 RMS 上升、RUL 下降，Spearman 可能是 $-0.85$。因此 ranking 使用：
\[
|\rho_{\mathrm{Spearman}}|
\]

\subsection{退化行为分数}

\term{退化行为分数}不是数据集官方给出的标签，也不是模型训练的损失函数。它是项目在特征分析阶段使用的一组\term{健康指标候选特征评价指标}，目的是回答：
\[
\text{这个特征本身的时间曲线，像不像一个能描述轴承退化过程的 HI？}
\]

它和 RUL 相关性有联系，但不是同一件事。RUL 相关性看的是“特征值和 RUL 标签是否相关”；退化行为分数看的是“特征沿时间演化时是否单调、平滑、跨轴承趋势一致，以及临近失效时是否收敛到相近水平”。因此，一个特征即使与 RUL 有相关性，也可能因为曲线噪声很大或跨轴承不稳定而不适合作为可解释的健康指标。

设某个特征在第 $b$ 个轴承上的时间序列为：
\[
x^{(b)}_1,x^{(b)}_2,\ldots,x^{(b)}_{T_b}
\]
其中 $T_b$ 是该轴承的样本数，样本已经按照 \code{timestep} 从早到晚排序。项目当前计算以下四个分数。

\begin{longtable}{L{3.8cm}L{10.5cm}}
\toprule
指标 & 解释 \\
\midrule
\endhead
\code{monotonicity} & 单调性。看同一轴承内特征是否主要朝一个方向变化。高分说明曲线大多持续上升或持续下降。 \\
\code{robustness} & 鲁棒性。看特征曲线是否平滑，是否不容易被局部噪声支配。高分说明曲线中的有效趋势强于短时抖动。 \\
\code{trendability} & 趋势一致性。看不同轴承的特征曲线形状是否相似。高分说明该特征在多根轴承上表现出相似退化模式。 \\
\code{prognosability} & 可预测性/末端一致性。看不同轴承临近失效时的特征值是否集中。高分说明该特征可能存在相对稳定的失效末端水平。 \\
\bottomrule
\end{longtable}

\subsubsection{Monotonicity：单调性}

单调性衡量一个特征在同一轴承内部是否主要沿一个方向变化。项目先计算相邻时间点差分：
\[
\Delta_t=x_{t+1}-x_t
\]
然后统计上升次数 $n_+$ 和下降次数 $n_-$：
\[
n_+=\#\{\Delta_t>0\}, \qquad n_-=\#\{\Delta_t<0\}
\]
单个轴承的单调性为：
\[
M_b=\frac{|n_+-n_-|}{T_b-1}
\]
多根轴承取平均：
\[
M=\frac{1}{B}\sum_{b=1}^{B}M_b
\]

如果一个特征几乎一直上升，$n_+$ 远大于 $n_-$，分数接近 1；如果它一会儿上升一会儿下降，$n_+$ 和 $n_-$ 接近，分数接近 0。

\subsubsection{Robustness：鲁棒性}

鲁棒性衡量特征曲线中的长期趋势是否强于局部噪声。项目用长度为 3 的滚动均值作为平滑曲线：
\[
\tilde{x}_t=\mathrm{rolling\_mean}_3(x_t)
\]
残差为：
\[
r_t=x_t-\tilde{x}_t
\]
单个轴承的鲁棒性为：
\[
R_b=\frac{1}{1+\frac{\mathrm{std}(r)}{\mathrm{std}(x)+\epsilon}}
\]
其中 $\epsilon=10^{-12}$ 用来避免除零。多根轴承取平均：
\[
R=\frac{1}{B}\sum_{b=1}^{B}R_b
\]

如果残差波动远小于原始特征自身的变化，说明曲线主要由平滑趋势构成，分数较高；如果曲线几乎全是短时抖动，残差很大，分数较低。

\subsubsection{Trendability：趋势一致性}

趋势一致性衡量同一个特征在不同轴承上的退化形状是否相似。由于不同轴承寿命长度不同，项目先把每根轴承的时间轴归一化到 $[0,1]$，再插值到相同数量的点：
\[
y^{(b)}=\mathrm{interp}\left(x^{(b)},K\right)
\]
其中 $K$ 是配置里的 \code{interpolation\_points}，默认 100。然后对所有轴承两两计算 Pearson 相关系数，并取绝对值平均：
\[
T=\frac{2}{B(B-1)}\sum_{i<j}\left|\rho\left(y^{(i)},y^{(j)}\right)\right|
\]

高分说明该特征在不同轴承上呈现相似曲线形状。需要注意：当前实现使用绝对相关系数，因此它强调“曲线形状强相关”，而不是严格要求所有轴承都同向变化。若后续论文特别强调方向一致性，可以在实验设置里说明这一点，或改成不取绝对值的版本。

\subsubsection{Prognosability：可预测性/末端一致性}

\code{prognosability} 在很多 HI 评价语境中可以理解为\term{可预测性}，但它不是模型预测准确率，而是特征曲线自身的一个性质。它关注的是：
\[
\text{不同轴承临近失效时，这个特征是否收敛到相似的末端水平？}
\]

项目中，对每根轴承取最后一个时间点的特征值：
\[
e_b=x^{(b)}_{T_b}
\]
同时计算该轴承全寿命内该特征的变化范围：
\[
a_b=\max_t x^{(b)}_t-\min_t x^{(b)}_t
\]
然后计算：
\[
P=\exp\left(
-\frac{\mathrm{std}(e_1,e_2,\ldots,e_B)}
{\mathrm{mean}(a_1,a_2,\ldots,a_B)+\epsilon}
\right)
\]

这个公式可以这样理解：
\begin{itemize}
  \item 分子 $\mathrm{std}(e_b)$ 表示不同轴承末端特征值有多分散。越分散，说明失效末端不稳定。
  \item 分母 $\mathrm{mean}(a_b)$ 表示该特征在单根轴承全寿命内本来有多大变化幅度。用它归一化后，可以避免只因为量纲不同而改变分数。
  \item 外层指数函数把结果压到 $(0,1]$。越接近 1，末端越一致；越接近 0，末端越分散。
\end{itemize}

例如，某个特征在三根轴承失效末端的值分别为 9、10、11，而各自全寿命变化范围约为 20、18、22。此时：
\[
\mathrm{std}(e)\approx0.82,\qquad \mathrm{mean}(a)\approx20
\]
\[
P\approx \exp(-0.82/20)\approx0.96
\]
这说明不同轴承接近失效时该特征收敛到相近水平，末端一致性较好。

相反，如果另一个特征的末端值为 5、30、80，变化范围约为 20、25、30，则：
\[
\mathrm{std}(e)\approx31.18,\qquad \mathrm{mean}(a)\approx25
\]
\[
P\approx \exp(-31.18/25)\approx0.29
\]
这说明不同轴承的末端值差异很大，很难把它解释为一个稳定的失效阈值。

\subsubsection{一个直观 case}

假设比较两个候选特征：\code{rms} 和一个随机噪声特征。

\begin{longtable}{L{3.2cm}L{5.5cm}L{5.5cm}}
\toprule
维度 & \code{rms} 类特征 & 随机噪声特征 \\
\midrule
\endhead
单调性 & 随磨损加剧逐渐上升，通常较高。 & 上下乱跳，通常较低。 \\
鲁棒性 & 有整体趋势，局部噪声相对较小。 & 主要由噪声构成，残差较大。 \\
趋势一致性 & 多根轴承可能都呈现早期平稳、后期上升。 & 不同轴承之间没有稳定形状。 \\
末端一致性 & 如果失效前振动能量都达到相近水平，则较高。 & 末端值随机，通常较低。 \\
\bottomrule
\end{longtable}

因此，\code{rms} 可能被推荐给 RUL 或健康状态任务；随机噪声即使偶然在某个轴承上和 RUL 有相关，也不应被过度解释为稳定退化特征。

\subsubsection{与 RUL 相关性的关系}

退化行为分数和 RUL 相关性互补：
\begin{itemize}
  \item \term{RUL 相关性}回答：这个特征和 RUL 标签在统计上是否相关？
  \item \term{退化行为分数}回答：这个特征的时间曲线是否符合退化过程应有的形态？
\end{itemize}

项目当前 RUL score 综合使用 Spearman 相关性、单调性、鲁棒性和趋势一致性：
\[
\text{rul\_score}
=0.35|\rho_s|
 +0.25\text{monotonicity}
 +0.20\text{robustness}
 +0.20\text{trendability}
\]
\code{prognosability} 当前写入 \code{degradation_scores.csv}，主要用于解释和诊断特征质量；它暂未进入当前 \code{rul\_score} 的加权公式。

\subsubsection{解释边界}

退化行为分数需要谨慎解读：
\begin{itemize}
  \item 它不是官方标签，也不是某个数据集给出的 ground truth。
  \item 它依赖同一轴承内部样本的时间顺序，因此不能在打乱时间顺序后计算。
  \item 对于非完整 run-to-failure 数据，\code{prognosability} 的“末端值”不一定是真正失效末端，需要谨慎解释。
  \item 高分只能说明该特征更像健康指标候选，不等于模型一定能取得最好预测精度。
\end{itemize}

\section{健康状态特征分析指标}

健康状态关注：\term{这个特征能不能区分 healthy / slight / moderate / severe？}

\begin{longtable}{L{4cm}L{10.2cm}}
\toprule
指标 & 解释 \\
\midrule
\endhead
\code{mutual_information} & 互信息。衡量特征和健康状态之间的信息量，越高说明越有关联。 \\
\code{fisher_score} & 类间差异与类内波动的比值。类间越分开、类内越集中，分数越高。 \\
\code{anova_f} & 方差分析 F 值。判断不同健康状态的均值是否明显不同。 \\
\code{state_0_mean} 等 & 每个健康状态下该特征的均值。可以看状态是否逐渐上移或下移。 \\
\code{state_0_std} 等 & 每个健康状态下该特征的标准差。可以看每个状态内部是否稳定。 \\
\bottomrule
\end{longtable}

项目当前 health score：
\[
\text{health\_score}
=0.4\cdot\mathrm{MI}
 +0.4\cdot\mathrm{Fisher}
 +0.2\cdot\mathrm{ANOVA}
\]
实际 ranking 前会对不同指标做归一化，避免量纲不一致。

\section{早期故障检测指标}

早期故障检测关注：\term{FPT 前后是否已经能分开？}

\begin{longtable}{L{4cm}L{10.2cm}}
\toprule
指标 & 解释 \\
\midrule
\endhead
\code{healthy_mean} & FPT 前样本的特征均值。 \\
\code{fault_mean} & FPT 后样本的特征均值。 \\
\code{mean_shift} & FPT 后均值减去 FPT 前均值。 \\
\code{healthy_std} & FPT 前波动。健康期波动太大，说明容易误报。 \\
\code{fault_std} & FPT 后波动。异常期波动太大，说明状态不稳定。 \\
\code{cohens_d} & 两组均值差除以合并标准差，表示效应量。绝对值越大，两组越分开。 \\
\code{auc} / \code{auc_abs} & ROC-AUC 的方向无关版本，$\max(\mathrm{AUC},1-\mathrm{AUC})$。越接近 1 越容易区分。 \\
\bottomrule
\end{longtable}

项目当前 early fault score：
\[
\text{early\_fault\_score}
=0.5|\mathrm{Cohen's}\ d|
 +0.5\cdot \mathrm{AUC}_{\mathrm{effective}}
\]

\section{故障类型特征分析}

当前阶段的 fault type 分析是弱版本，因为还没有加入 BPFO/BPFI/BSF/FTF 等轴承机理频率特征。现阶段只能用 manual/tsfresh 统计特征去看它们是否能弱区分：
\[
\text{normal},\quad \text{degraded\_unknown},\quad \text{outer/inner/cage/compound}
\]

这可以作为探索，但不能过度解释为“已经识别了具体故障机理”。真正的故障机理解释后续应加入：
\begin{itemize}
  \item BPFO energy：外圈故障特征频率能量。
  \item BPFI energy：内圈故障特征频率能量。
  \item BSF energy：滚动体故障特征频率能量。
  \item FTF energy：保持架故障特征频率能量。
  \item envelope spectrum peak：包络谱峰值。
  \item harmonic/sideband energy：谐波和边带能量。
\end{itemize}

\section{哪些特征通常适合哪些任务}

\begin{longtable}{L{4cm}L{4.2cm}L{6cm}}
\toprule
特征族 & 常适合任务 & 原因和注意事项 \\
\midrule
\endhead
\code{rms}, \code{std}, \code{mean_abs} & RUL、HealthState & 代表整体振动能量，退化时通常稳定上升。若被用作 HI source，需要标注泄漏风险。 \\
\code{ptp} & HealthState、EarlyFault & 对峰值跨度敏感，严重退化和冲击会增大，但可能受异常点影响。 \\
\code{kurtosis} & EarlyFault & 对局部冲击敏感，适合发现早期缺陷，但不一定单调，不一定适合 RUL 主特征。 \\
\code{crest_factor} & EarlyFault & 看尖峰相对 RMS 是否突出，适合突发冲击检测。 \\
\code{impulse_factor} & EarlyFault & 看尖峰相对平均绝对幅值是否突出，噪声敏感。 \\
\code{clearance_factor} & EarlyFault & 轴承故障中常用的冲击敏感指标，解释时要结合曲线。 \\
\code{spectral_centroid} & RUL、HealthState & 反映频谱能量中心迁移，但可能受工况变化影响。 \\
\code{rms_frequency} & RUL、HealthState & 反映整体频率结构变化。 \\
\code{spectral_entropy} & HealthState、EarlyFault & 反映频谱复杂度或宽频激励，解释性不如 RMS 直观。 \\
\code{tsfresh_abs_energy} 等能量特征 & RUL & 能量类特征通常与退化相关，但需要跨轴承验证。 \\
\code{tsfresh_kurtosis} 等冲击特征 & EarlyFault & 对尖峰和冲击敏感，但噪声风险高。 \\
\bottomrule
\end{longtable}

\section{如何阅读分析结果}

建议按以下顺序看：
\begin{enumerate}
  \item \term{看 \code{analysis_report.json}}：确认 \code{ok=true}、\code{fit_scope=train_only}、是否有 leakage warning、PHM2012 是否跳过 fault type。
  \item \term{看 \code{feature_summary.csv}}：先排除 NaN 多、Inf 多、常数列、方差极小的特征。
  \item \term{看 \code{feature_ranking.csv}}：分别按 \code{rank_rul}、\code{rank_health}、\code{rank_early_fault} 排序。
  \item \term{看 \code{feature_cards.csv}}：理解每个特征的物理意义、适合任务、原因和 caveat。
  \item \term{看 \code{feature_recommendations.md}}：快速获得每类任务推荐特征。
  \item \term{看 \code{figures/}}：用曲线和箱线图人工确认，不要只相信分数。
\end{enumerate}

\section{常见误解澄清}

\begin{longtable}{L{5cm}L{9.2cm}}
\toprule
误解 & 正确理解 \\
\midrule
\endhead
“加速度值”和“振动值”是不是两个东西？ & 在本项目语境里，原始 CSV 的最后两列是振动加速度信号。口语里说振动值，通常就是指这些加速度采样值。 \\
RUL 是真实标签吗？ & 对 run-to-failure 数据，可根据最终失效时间构造 RUL；但 piecewise RUL 依赖 FPT，是一种任务标签设计。 \\
FPT 是官方标签吗？ & 不是。当前 FPT 是项目用 HI 和 3$\sigma$ 规则计算出的退化起点。 \\
healthy / slight / moderate / severe 是原始数据给的吗？ & 不是逐样本官方标签，而是项目基于 FPT 和退化进度构造的健康状态标签。 \\
early fault 是不是严重故障？ & 不是。当前 early fault 表示已经过 FPT，进入可观测异常阶段。 \\
RMS 排名高就证明 RMS 独立发现故障吗？ & 不一定。如果 RMS 被用来构造 HI/FPT，就必须说明 label-source 风险。 \\
test set 能不能用来挑特征？ & 不应该。特征 ranking 应使用 train split，否则评估不可信。 \\
PHM2012 能不能做 fault type？ & 当前不主打，因为缺少 XJTU-SY 那种轴承级 fault element 标签。 \\
\bottomrule
\end{longtable}

\section{本项目命令示例}

XJTU-SY 上运行完整特征分析：
\begin{verbatim}
uv run bp --config-name smoke \
  mode=analyze_features \
  dataset=xjtu_sy \
  split=xjtu_leave_one_bearing_out \
  feature=manual_basic \
  label=degradation_basic \
  analysis=full_feature_analysis \
  dataset.root=/Users/nev8r/Desktop/phm2/data/loader_roots/xjtu \
  split.condition_id=35Hz12kN \
  split.test_bearing_id=Bearing1_5 \
  split.val_bearing_id=Bearing1_4
\end{verbatim}

运行后重点查看：
\begin{verbatim}
artifacts/runs/<run_id>/analysis/feature_ranking.csv
artifacts/runs/<run_id>/analysis/feature_cards.csv
artifacts/runs/<run_id>/analysis/feature_recommendations.md
artifacts/runs/<run_id>/analysis/figures/
\end{verbatim}

\section{缩写表}

\begin{longtable}{L{3cm}L{5cm}L{6cm}}
\toprule
缩写 & 全称 & 中文解释 \\
\midrule
\endhead
PHM & Prognostics and Health Management & 故障预测与健康管理 \\
RUL & Remaining Useful Life & 剩余使用寿命 \\
HI & Health Indicator & 健康指标 \\
FPT & First Predictable Time & 首次可预测时间/首次明显退化点 \\
FFT & Fast Fourier Transform & 快速傅里叶变换 \\
rFFT & Real-input FFT & 实数输入的一侧 FFT \\
RMS & Root Mean Square & 均方根 \\
STD & Standard Deviation & 标准差 \\
PTP & Peak-to-Peak & 峰峰值 \\
AUC & Area Under ROC Curve & ROC 曲线下面积 \\
MI & Mutual Information & 互信息 \\
ANOVA & Analysis of Variance & 方差分析 \\
BPFO & Ball Pass Frequency Outer Race & 外圈故障特征频率 \\
BPFI & Ball Pass Frequency Inner Race & 内圈故障特征频率 \\
BSF & Ball Spin Frequency & 滚动体自转故障频率 \\
FTF & Fundamental Train Frequency & 保持架故障特征频率 \\
\bottomrule
\end{longtable}

\section{一句话总结}

本项目的核心思想是：把轴承全寿命振动信号转换成可解释的时间域和频域特征，再用 HI 和 FPT 构造 RUL、健康状态和早期故障标签，最后通过 train-only 的特征分析判断哪些特征适合寿命预测、哪些适合状态评估、哪些适合早期报警，并明确标注标签来源风险，避免把由标签构造带来的优势误当成独立发现能力。

\end{document}
